\section{Introduction}\label{sec:intro}

Motor insurers handle claims in large numbers, and much of the cost lies in manual assessment: surveyors inspect photographs, read the claimant's account and check the policy record, and different assessors often reach different conclusions about the same case~\cite{ouf2026,cheng2026}. Insurers therefore automate digital first notice of loss, image-based damage detection, repair-cost estimation and fraud screening~\cite{arconzo2026,subhash2026,ouf2026}; in one deployment, conversational agents with damage detection cut processing time from 30 to 12 minutes~\cite{arconzo2026}. The risks are asymmetric. Approving a fraudulent claim loses money, auditing a legitimate one delays payment and adds handling cost, and a small, imbalanced fraud class makes the trade-off harder to control~\cite{bekkaye2026}. Deciding which claims to settle automatically, which to audit and which to give to a person is a practical problem worth studying.

This paper studies that routing decision, which we call \emph{claim gating}. Each claim has a photograph of the damaged vehicle, a short categorical description and eleven numerical fields, and the system assigns \emph{Auto-Approve}, \emph{Reject/Fraud-Audit} or \emph{Escalate-Human}. The task is hard because the photograph may show more or less damage than reported or be blurred, dark or cropped, the description may be incomplete, and errors carry different costs. We study it offline on a constructed 1,000-claim benchmark pairing CarDD photographs~\cite{wang2023cardd} with unrelated records from a public claims data set~\cite{kaggleclaims}; its labels follow documented construction rules, so results describe routing under those rules, not on real claims. Repair-cost estimation and forgery detection are out of scope.

The 2026 studies we surveyed leave four gaps. First, damage assessment and fraud prediction are studied separately~\cite{subhash2026}: vision models judge the photograph alone~\cite{ouf2026,li2026part,dalal2026,cheng2026}, and fraud models score the claim table without it~\cite{bekkaye2026,ebinezer2026,vandervorst2026}. Second, outputs are mostly binary; routing uncertain cases to people has been developed for payment and healthcare fraud~\cite{liu2026daunt,sourav2026,dadi2026,syed2026}, not for motor claims with photographs. Third, calibration and expected cost have been analysed for tabular motor-fraud models~\cite{bekkaye2026} and abstention for image-only damage assessment~\cite{shao2026}, but no study checks whether the photograph agrees with the claimed severity before deciding. Fourth, motor-claim systems that use language models generate text inside multi-stage pipelines~\cite{arconzo2026,mitrenga2026,phuengpanyaloet2026,hatzesberger2026}; none reports the per-call latency and cost of a routing decision or compares a closed-set decision model with a generative model asked the same questions, although per-path latency has been reported for transaction fraud~\cite{sourav2026}. What is missing is a calibrated, cost-aware, three-way routing policy that uses all three kinds of evidence and checks them against one another.

We build such a router on Gemini Decision~1, a hosted multimodal model that scores closed-set questions in one forward pass and returns option probabilities instead of text~\cite{google2026decisions}. One request carries the photograph, the description and the numbers as labelled text, and asks a three-option routing question and four yes/no checks: whether the visible damage matches the reported severity, whether the photograph is clear, whether the details are complete, and whether the amounts match the severity. Two checks compare modalities directly (image with text, numbers with text). Because the model generates no text and has a lower list price, a decision may be cheaper and faster than a generative one; we measure this against a generative model given the same questions, resolution and combiner, which compares two hosted products rather than two mechanisms.

A 24-parameter logistic combiner maps the answers to route probabilities, and a claim is approved or audited automatically only when the routing confidence clears a threshold chosen on validation to minimise a misrouting cost in which a wrong approval is the most expensive error; otherwise it is escalated, following the reject option~\cite{chow1970} and selective classification~\cite{geifman2017selective}. We assess accuracy, calibration~\cite{guo2017calibration,naeini2015,brier1950}, risk--coverage, misrouting and handling cost, latency and cost, on a fixed split and in repeated cross-validation, against single-modality, trained-fusion and generative baselines, with the same shortcut controls for every method.

The contributions of this paper are as follows.
\begin{itemize}
\item \textbf{A documented tri-modal claim-gating benchmark.} It contains 1,000 claim--photograph pairs with an image, a description, eleven numerical fields, three routing labels and a fixed split. We document its construction rules and measure the shortcuts they create; the record set will be released subject to the terms of both source data sets, and the images remain subject to the CarDD licence.
\item \textbf{A cost-aware three-way gating formulation.} Claim gating is formalised as selective three-way classification with an escalation class and an asymmetric cost matrix. At moderate cost weights the gate lowers misrouting cost against escalating every claim in cross-validation; once every human review is charged, it is no better than escalating everything.
\item \textbf{A single-request closed-set decision design.} Five questions, including image--text and number--text consistency checks, are answered in one request and combined by 24 parameters. The answers are more accurate routing features than a generative model's answers to the same questions, and, because the questions were written from the routing rules, few labels are needed: with 25 labelled claims the router reaches 0.825 accuracy, against 0.754 for a learner given the rules as features and 0.609 for trained fusion.
\item \textbf{An evaluation protocol with released code and cached responses.} The configuration was frozen before it was scored on test, and the router is compared with text, numerical, image, early-fusion, generative and rule-aware baselines using exact tests, corrected cross-validation tests, learning curves, a stability re-test, a photograph-degradation control and ablations. The benchmark build files were not kept, so the benchmark cannot be rebuilt from the source data sets.
\end{itemize}
